\section{Entropies of flat densities and harmonic weights}

\begin{lemma}[Entropy of a flat density operator]%
    \label{thm:peps_flat_density_entropy}
    \lean{Matrix.IsHermitian.eigenvalues_eq_zero_or_inv_of_mul_self_eq_inv_smul,
        vonNeumannEntropy_of_mul_self_eq_inv_smul}
    \leanok
    \uses{def:von_neumann_entropy}
    Let $\rho$ be a finite-dimensional positive semidefinite matrix of
    trace one. As in the spectral step of the proof of
    \cite[Theorem~6.9]{Schuch2010PEPS}, if $\rho^2=r^{-1}\rho$ for a positive
    real number $r$, then $S(\rho)=\log r$.
\end{lemma}

\begin{proof}\leanok
    For each eigenvalue $\lambda$, the identity gives
    $\lambda(\lambda-r^{-1})=0$. In either case
    $-\lambda\log\lambda=\lambda\log r$, with the value at zero
    interpreted continuously. Summing and using
    $\sum\lambda=\operatorname{tr}\rho=1$ proves $S(\rho)=\log r$.
\end{proof}

\begin{lemma}[All finite R\'enyi orders of a flat density]%
    \label{thm:peps_flat_density_renyi_entropy}
    \lean{trace_cfc_rpow_of_mul_self_eq_inv_smul,
        renyiEntropy_of_mul_self_eq_inv_smul,
        renyiEntropy_of_mul_self_eq_inv_smul_of_rank}
    \leanok
    \uses{def:renyi_entropy}
    Let $\rho\geq0$ have trace one and satisfy $\rho^2=r^{-1}\rho$, with
    $r>0$. For every $\alpha>0$,
    $\operatorname{tr}(\rho^\alpha)=r^{1-\alpha}$ and
    $S_\alpha(\rho)=\log r$. If $\operatorname{rank}\rho=r$, the same
    entropy formula also holds at order zero.
\end{lemma}

\begin{proof}\leanok
    \uses{thm:peps_flat_density_entropy}
    Every eigenvalue is zero or $r^{-1}$. For positive order,
    $\lambda^\alpha=\lambda r^{1-\alpha}$. Summing uses trace one;
    taking logarithms proves the entropy formula away from order one.
    The von Neumann and rank conventions give the remaining orders.
\end{proof}

\begin{lemma}[Harmonic weights inside the tail radius]
    \label{lem:harmonic_tail_weights}
    \lean{Entropy.harmonic_weight_le_half,
        Entropy.harmonic_weight_ratio_le,
        Entropy.le_tailRadius_of_le}
    \leanok
    Let $Z\ge4$ and $c\ge1$, and put $a=2/(Zc)$. Then $a\le1/2$ and
    $a/(1-a)\le4/(Zc)$. If moreover $\vartheta\ge0$,
    $1\le\mathcal B\le K'c$ and $Z\ge128\sqrt{(1+\vartheta)K'}$, then
    \begin{align}
        \frac4{Zc}\le\frac1{32\sqrt{(1+\vartheta)\mathcal B}},
        \notag
    \end{align}
    so the tilt $a/(1-a)$ lies inside the tail radius of the marginal tail
    bound. These are the weight estimates of the proof of
    \cite[Lemma~3.2]{OpenAI2026AreaLaw}.
\end{lemma}

\begin{proof}\leanok
    The first two bounds are elementary: $a\le1/2$ gives
    $a/(1-a)\le2a$. For the last, $\sqrt{(1+\vartheta)\mathcal B}\le
        \sqrt{(1+\vartheta)K'}\sqrt c\le\sqrt{(1+\vartheta)K'}\,c$, and
    multiplying the lower bound on $Z$ by $c$ gives the claim.
\end{proof}
